\documentclass[10pt,aspectratio=169]{beamer}
% Preámbulo modular para presentaciones Beamer (Modelación Estadística)
% Diseño y paleta institucional del Tecnológico de Monterrey

\documentclass[aspectratio=169,xcolor={dvipsnames,table}]{beamer}

\usepackage[utf8]{inputenc}
\usepackage[spanish,mexico]{babel}
\usepackage{amsmath,amssymb,amsfonts,mathtools}
\usepackage{graphicx}
\usepackage{listings}
\usepackage{booktabs}
\usepackage{tikz}
\usetikzlibrary{matrix,arrows,decorations.pathmorphing,shapes.geometric,positioning}

% Paleta Institucional Tecnológico de Monterrey
\definecolor{TecAzul}{HTML}{0039A6}       % Azul TEC: color primario institucional
\definecolor{TecAzulOscuro}{HTML}{1A2E51} % Azul marino oscuro
\definecolor{TecRojo}{HTML}{EC2661}       % Rojo de acento (paleta miTec)
\definecolor{TecGris}{HTML}{646464}       % Gris neutro para comentarios y subtítulos
\definecolor{TecFondoBloque}{HTML}{F4F6F9}% Fondo suave para bloques

% Configuración de Tema Beamer
\usetheme{Boadilla}
\usecolortheme[named=TecAzulOscuro]{structure}
\setbeamercolor{alerted text}{fg=TecRojo}
\setbeamercolor{palette primary}{bg=TecAzulOscuro,fg=white}
\setbeamercolor{palette secondary}{bg=TecAzul,fg=white}
\setbeamercolor{palette tertiary}{bg=TecAzulOscuro!80!black,fg=white}
\setbeamercolor{title}{fg=TecAzulOscuro,bg=TecFondoBloque}
\setbeamercolor{frametitle}{fg=TecAzulOscuro,bg=TecFondoBloque}

% Bloques personalizados elegantes
\setbeamercolor{block title}{bg=TecAzulOscuro,fg=white}
\setbeamercolor{block body}{bg=TecFondoBloque,fg=black}
\setbeamercolor{block title alerted}{bg=TecRojo,fg=white}
\setbeamercolor{block body alerted}{bg=TecRojo!8,fg=black}
\setbeamercolor{block title example}{bg=TecAzul,fg=white}
\setbeamercolor{block body example}{bg=TecAzul!8,fg=black}

% Redondear bordes y añadir sombra sutil
\setbeamertemplate{blocks}[rounded][shadow=true]
\setbeamertemplate{navigation symbols}{}

% Pie de página limpio con número de diapositiva
\setbeamertemplate{footline}{
  \leavevmode%
  \hbox{%
  \begin{beamercolorbox}[wd=.7\paperwidth,ht=2.25ex,dp=1ex,left]{author in head/foot}%
    \hspace*{2ex}\insertshortauthor~~(\insertshortinstitute) \hfill \insertshorttitle
  \end{beamercolorbox}%
  \begin{beamercolorbox}[wd=.3\paperwidth,ht=2.25ex,dp=1ex,right]{date in head/foot}%
    \insertshortdate{}\hspace*{2em}
    \insertframenumber{} / \inserttotalframenumber\hspace*{2ex}
  \end{beamercolorbox}}%
  \vskip0pt%
}

% Configuración de Listings de Python para visualización pedagógica de código
\lstset{
    language=Python,
    basicstyle=\ttfamily\footnotesize,
    keywordstyle=\color{TecAzulOscuro}\bfseries,
    stringstyle=\color{TecRojo},
    commentstyle=\color{TecGris}\itshape,
    showstringspaces=false,
    breaklines=true,
    frame=lines,
    rulecolor=\color{TecAzulOscuro!30},
    backgroundcolor=\color{TecFondoBloque},
    aboveskip=0.5em,
    belowskip=0.5em,
    literate={á}{{\'a}}1 {é}{{\'e}}1 {í}{{\'i}}1 {ó}{{\'o}}1 {ú}{{\'u}}1
             {Á}{{\'A}}1 {É}{{\'E}}1 {Í}{{\'I}}1 {Ó}{{\'O}}1 {Ú}{{\'U}}1
             {ñ}{{\~n}}1 {Ñ}{{\~N}}1 {¿}{{?`}}1 {¡}{{!`}}1
}

% Metadatos institucionales por defecto
\title[Modelación Estadística]{Modelación Estadística para Ciencia de Datos e Ingeniería}
\author[J. Castillo Colmenares]{Juliho David Castillo Colmenares}
\institute[Tec de Monterrey]{Tecnológico de Monterrey \\ Escuela de Ingeniería y Ciencias}
\date{\today}

% Comandos y macros estadísticas compartidas para las presentaciones Beamer (Metropolis)

% Conjuntos numéricos básicos
\newcommand{\R}{\mathbb{R}}
\newcommand{\N}{\mathbb{N}}
\newcommand{\Q}{\mathbb{Q}}
\newcommand{\Z}{\mathbb{Z}}
\newcommand{\set}[1]{\left\{ #1 \right\}}
\newcommand{\abs}[1]{\left|#1\right|}
\newcommand{\norm}[1]{\left\| #1 \right\|}

% Comandos estadísticos y probabilísticos del curso
\newcommand{\Var}{\operatorname{Var}}
\newcommand{\cov}{\operatorname{Cov}}
\newcommand{\corr}{\rho}
\newcommand{\comb}[2]{\begin{pmatrix} #1 \\ #2 \end{pmatrix}}
\newcommand{\s}{\sigma}
\newcommand{\particion}{\mathfrak{P}}
\newcommand{\card}[1]{\#\left( #1 \right)}
\newcommand{\seq}[3]{\set{#1}_{#2}^{#3}}
\newcommand{\E}{\mathbb{E}}
\newcommand{\Prob}{\mathbb{P}}

% Entornos pedagógicos y cajas en español académico natural
\newenvironment{discusion}{%
  \begin{alertblock}{Pregunta de Discusión}%
}{%
  \end{alertblock}%
}

\newenvironment{reflexion}{%
  \begin{alertblock}{Reflexión para el Aula}%
}{%
  \end{alertblock}%
}

\newenvironment{paradoja}{%
  \begin{alertblock}{Paradoja de Exploración}%
}{%
  \end{alertblock}%
}

\newenvironment{motivacion}{%
  \begin{exampleblock}{Motivación: Ciencia de Datos en la Práctica}%
}{%
  \end{exampleblock}%
}

\newenvironment{retoclase}{%
  \begin{block}{Reto Rápido en Equipo}%
}{%
  \end{block}%
}


\title[Transformación de variables]{Sección 5.1: Transformación de variables}
\subtitle{Distribuciones de $Y=g(X)$}
\author[J. Castillo Colmenares]{Juliho Castillo Colmenares}
\institute{Tecnológico de Monterrey}
\date{}

\begin{document}
\begin{frame}[plain]\vspace{-0.8cm}\titlepage\end{frame}

\begin{frame}{Hoja de ruta --- capítulo 5}
  \footnotesize
  \begin{itemize}\itemsep=0.08cm
    \item \textbf{05.00} Introducción inferencial
    \item \textbf{05.01} Transformación de variables \emph{(hoy)}
    \item \textbf{05.02} Distribuciones de funciones de variables aleatorias
    \item \textbf{05.03--05.05} Distribuciones $\chi^2$, $t$ y $F$
  \end{itemize}
  \begin{block}{Objetivo didáctico}
    Obtener la densidad de $Y=g(X)$ mediante transformaciones lineales,
    inversas y jacobianos, incluyendo el caso no monótono.
  \end{block}
\end{frame}

\begin{frame}{Motivación: cambiar la escala}
  \footnotesize
  Si $X$ tiene una distribución conocida y $Y=g(X)$, necesitamos describir la
  distribución de $Y$ para poder analizar datos transformados.
  \begin{itemize}\itemsep=0.12cm
    \item Una transformación puede normalizar datos.
    \item También puede estabilizar varianzas o comprimir valores extremos.
    \item La transformación debe conservar la probabilidad total.
  \end{itemize}
\end{frame}

\begin{frame}{Transformación afín}
  \footnotesize
  Si $Y=aX+b$, $a\neq0$, y $X$ tiene media $\mu_X$ y varianza $\sigma_X^2$,
  \[
    \mu_Y=a\mu_X+b,\qquad \sigma_Y^2=a^2\sigma_X^2.
  \]
  \begin{alertblock}{Caso normal}
    Si $X\sim N(\mu,\sigma^2)$, entonces
    $Y\sim N(a\mu+b,a^2\sigma^2)$.
  \end{alertblock}
\end{frame}

\begin{frame}{Cambio de variable monótono}
  \footnotesize
  Para $Y=g(X)$ con $g$ estrictamente monótona y diferenciable:
  \[
    f_Y(y)=f_X\!\left(g^{-1}(y)\right)
    \left|\frac{d}{dy}g^{-1}(y)\right|.
  \]
  \begin{block}{Idea de la derivación}
    Se expresa primero $F_Y(y)=F_X(g^{-1}(y))$ y después se deriva usando la
    regla de la cadena.
  \end{block}
\end{frame}

\begin{frame}{Jacobiano de la inversa}
  \footnotesize
  Si $x=g^{-1}(y)$ y $J(y)=dx/dy$, la fórmula se escribe
  \[
    f_Y(y)=f_X(x)|J(y)|.
  \]
  \begin{exampleblock}{Lectura geométrica}
    El valor absoluto del jacobiano corrige el cambio de escala entre los
    intervalos de $x$ y de $y$.
  \end{exampleblock}
\end{frame}

\begin{frame}{Caso no monótono}
  \footnotesize
  Si varios valores de $x$ producen el mismo $y$, se suman todas las
  preimágenes:
  \[
    f_Y(y)=\sum_{x:g(x)=y}f_X(x)\left|\frac{dx}{dy}\right|.
  \]
  \begin{alertblock}{Diferencia clave}
    No basta con escoger una sola inversa: cada rama aporta probabilidad.
  \end{alertblock}
\end{frame}

\begin{frame}{Transformaciones comunes}
  \footnotesize
  La nota enumera transformaciones útiles para datos positivos o sesgados:
  \begin{itemize}\itemsep=0.11cm
    \item $Y=\ln X$: comprime valores grandes.
    \item $Y=\sqrt X$: suaviza la asimetría de conteos.
    \item $Y=X^{1/3}$: alternativa para asimetría pronunciada.
    \item Box--Cox: $Y=(X^\lambda-1)/\lambda$.
  \end{itemize}
\end{frame}

\begin{frame}{Aplicación exacta 1/4: transformación normal}
  \footnotesize
  Para $X\sim N(\mu,\sigma^2)$ y $Y=aX+b$:
  \[
    Y\sim N(a\mu+b,a^2\sigma^2).
  \]
  \begin{block}{Interpretación}
    Una transformación lineal desplaza la media y escala la varianza, pero
    conserva la familia normal.
  \end{block}
\end{frame}

\begin{frame}{Aplicación exacta 1/4: resolución}
  \footnotesize
  \begin{enumerate}\itemsep=0.13cm
    \item Identificar $a$ y $b$.
    \item Transformar la media: $a\mu+b$.
    \item Transformar la varianza: $a^2\sigma^2$.
  \end{enumerate}
  Este es el ejemplo de reproductividad de la distribución normal de la nota.
\end{frame}

\begin{frame}{Aplicación exacta 2/4: exponencial a Pareto}
  \footnotesize
  Sea $X\sim\operatorname{Exp}(1)$ y $Y=e^X$. Entonces $g^{-1}(y)=\ln y$ y
  $d(\ln y)/dy=1/y$.
  \[
    f_Y(y)=e^{-\ln y}\frac1y=\frac1{y^2},\qquad y>1.
  \]
\end{frame}

\begin{frame}{Aplicación exacta 2/4: resolución}
  \footnotesize
  \begin{block}{Soporte}
    Como $X>0$, la transformación creciente produce $Y>1$.
  \end{block}
  \begin{alertblock}{Resultado}
    La densidad $1/y^2$ para $y>1$ es una distribución de Pareto con
    parámetro $\alpha=1$.
  \end{alertblock}
\end{frame}

\begin{frame}{Aplicación exacta 3/4: transformación log-normal}
  \footnotesize
  Si $X\sim N(\mu,\sigma^2)$ y $Y=e^X$, la inversa es $\ln y$ y $y>0$.
  \[
    f_Y(y)=\frac{1}{\sigma\sqrt{2\pi}}
    \exp\!\left[-\frac{(\ln y-\mu)^2}{2\sigma^2}\right]\frac1y.
  \]
\end{frame}

\begin{frame}{Aplicación exacta 3/4: resolución}
  \footnotesize
  \begin{block}{Resultado de la nota}
    $Y$ sigue una distribución log-normal $LN(\mu,\sigma^2)$.
  \end{block}
  \begin{exampleblock}{Lectura}
    El factor $1/y$ es el jacobiano de la transformación inversa y evita
    confundir la densidad de $X$ con la de $Y$.
  \end{exampleblock}
\end{frame}

\begin{frame}{Aplicación exacta 4/4: seno no monótono}
  \footnotesize
  Para $X\sim U(0,2\pi)$ y $Y=\sin X$, cada $y\in(-1,1)$ tiene dos preimágenes.
  \[
    x_1=\arcsin y,\qquad x_2=\pi-\arcsin y.
  \]
  \begin{block}{Suma de ramas}
    Cada rama contribuye $1/(2\pi\sqrt{1-y^2})$.
  \end{block}
\end{frame}

\begin{frame}{Aplicación exacta 4/4: resolución}
  \footnotesize
  \[
    f_Y(y)=\frac{1}{\pi\sqrt{1-y^2}},\qquad -1<y<1.
  \]
  \begin{alertblock}{Resultado}
    Se obtiene la densidad del arcoseno: ignorar una de las dos preimágenes
    produciría la mitad de la probabilidad correcta.
  \end{alertblock}
\end{frame}

\begin{frame}{Puente computacional 1/4: soporte}
  \footnotesize
  La transformación debe llevar también el soporte:
  \begin{center}
    $X>0\xrightarrow{e^x}Y>1$,\qquad
    $X\in(0,2\pi)\xrightarrow{\sin}Y\in(-1,1)$.
  \end{center}
  El soporte es una comprobación rápida de cualquier densidad transformada.
\end{frame}

\begin{frame}{Puente computacional 2/4: jacobiano}
  \footnotesize
  \[
    J(y)=\frac{d}{dy}g^{-1}(y).
  \]
  \begin{itemize}\itemsep=0.12cm
    \item $g^{-1}(y)=\ln y\Rightarrow |J(y)|=1/y$.
    \item $g^{-1}(y)=\arcsin y\Rightarrow |J(y)|=1/\sqrt{1-y^2}$.
  \end{itemize}
\end{frame}

\begin{frame}{Puente computacional 3/4: conservación}
  \footnotesize
  En cada caso, la densidad transformada debe integrar uno sobre su soporte.
  \begin{block}{Chequeo de consistencia}
    El jacobiano redistribuye masa; no crea ni elimina probabilidad.
  \end{block}
  Las dos ramas del seno explican el factor $2$ de la densidad del arcoseno.
\end{frame}

\begin{frame}{Puente computacional 4/4: decisión de método}
  \footnotesize
  \begin{tabular}{ll}
    \hline
    Forma de $g$ & Acción \\
    \hline
    Afín & transformar media y varianza \\
    Monótona & una inversa y un jacobiano \\
    No monótona & sumar todas las ramas \\
    \hline
  \end{tabular}
  \begin{alertblock}{Regla práctica}
    La forma de la función determina la forma de la solución.
  \end{alertblock}
\end{frame}

\begin{frame}{Síntesis de la sección}
  \footnotesize
  \begin{itemize}\itemsep=0.12cm
    \item Las transformaciones afines preservan la normalidad.
    \item El cambio de variable usa la inversa y su jacobiano.
    \item Las funciones no monótonas requieren sumar preimágenes.
    \item El soporte y la integral total verifican el resultado.
  \end{itemize}
\end{frame}

\begin{frame}{Transición curricular}
  \footnotesize
  \begin{block}{Lo que queda establecido}
    Ya podemos obtener distribuciones de transformaciones de una variable.
  \end{block}
  \begin{alertblock}{Siguiente sección}
    \textbf{05.02 Distribuciones de probabilidad de funciones de variable
    aleatoria}: extensión a funciones de una o varias variables y superficies
    de nivel.
  \end{alertblock}
\end{frame}
\end{document}
